\section{引言}

目标围捕控制要求一群智能体围绕目标自组织，使得目标位于智能体生成的
凸包内，同时智能体之间避免碰撞。与传统的编队跟踪不同，最终形状并非
预先固定，而是由智能体与人工势场（APF）项之间的相互作用生成。
这一特性使目标围捕控制具有吸引力，但也使分析变得微妙，
因为同一个APF项既要塑造最终的凸包，又要防止碰撞。

对于低阶智能体，APF项的作用相对透明。在单积分器模型中，
可以直接施加形式为 $-\kappa_i p_{i0}+\phi_i$ 的速度，
使得中心关系由 $\sum_i\phi_i=0$ 保证。
在二阶模型中，可以直接将APF力放入相对加速度动力学中，
并构造一个同时包含动能和APF势能的能量函数。
由于APF力是势能的负梯度，力项可以抵消势能的导数，
并通过证明障碍势能保持有界来证明避碰。

对于异构高阶智能体，情况要困难得多。智能体可能具有不同的相对阶，
而APF仍然是位置级别的量。要在高阶控制器中使用它，
必须对其进行微分、积分、滤波，或将其嵌入递归反步构造中。
每种操作都会改变问题的解析结构。
对APF进行微分是脆弱的，因为邻域集合可能切换，且APF在安全边界附近是奇异的。
积分或滤波APF可以使信号平滑，但同时也会削弱障碍证明所需的奇异性。
因此，控制器可能仍然能证明凸包围捕，却无法证明全时避碰。

本笔记因此被写成尝试的记录，而非最终的完整控制器。
我们考察了四个自然的方向。
第一，基于积分的设计跟踪APF的重复积分。由于重复积分在所有智能体上仍然
求和为零，这给出了简单的凸包论证，但凸包可能无界增长，且避碰未被保证。
第二，APF滤波滑模设计避免了对原始APF的微分，并给出了滤波APF参考量的
有限时间跟踪，但滤波破坏了直接的障碍能量抵消。
第三，精确反步尝试展示了为何反复微分APF不是一条稳健的路线。
最后，指令滤波反步尝试减少了对APF导数的需求，
但留下了残余滤波/跟踪误差，必须通过额外的安全论证来处理。

这些尝试的目的在于精确地隔离出难点所在。
如果只需要目标位于智能体凸包内（可能伴随不断扩张的编队），
凸包围捕相对容易获得。
避碰是更困难的部分：它要求安全集合的前向不变性，
而不仅仅是有界的跟踪误差或发散的APF项。
讨论表明，完整的高阶围捕控制器很可能需要将高阶跟踪机制与
显式的安全性证明（如障碍李雅普诺夫函数或高阶控制障碍函数）结合起来。